\documentclass[11pt,a4paper]{article}

\usepackage[a4paper,margin=22mm]{geometry}
\usepackage[T1]{fontenc}
\usepackage[utf8]{inputenc}
\usepackage{lmodern}
\usepackage{microtype}
\usepackage{amsmath,amssymb,amsthm,mathtools,mathrsfs}
\usepackage{booktabs,array,tabularx}
\usepackage{enumitem}
\usepackage{xcolor}
\usepackage{hyperref}
\usepackage{fancyhdr}
\usepackage{titlesec}

\ifdefined\pdfinfoomitdate\pdfinfoomitdate=1\fi
\ifdefined\pdftrailerid\pdftrailerid{}\fi
\ifdefined\pdfsuppressptexinfo\pdfsuppressptexinfo=-1\fi

\definecolor{deepblue}{RGB}{28,65,112}

\hypersetup{
  colorlinks=true,
  linkcolor=deepblue,
  citecolor=deepblue,
  urlcolor=deepblue,
  pdftitle={A Layered Framework for Compact-Resolvent Spectral Encodings},
  pdfauthor={Stassis Stashkevichyus; Stassis Research Program},
  pdfsubject={Compact resolvent; heat trace; spectral measure; complex-time encoding; theta kernels; readout factorization},
  pdfkeywords={compact resolvent, heat trace, spectral measure, Laplace transform, theta function, spectral encoding}
}

\setlength{\headheight}{14pt}
\pagestyle{fancy}
\fancyhf{}
\fancyhead[C]{Compact-Resolvent Spectral Encodings}
\fancyfoot[C]{\thepage}
\renewcommand{\headrulewidth}{0.4pt}
\fancypagestyle{plain}{
  \fancyhf{}
  \fancyhead[C]{Compact-Resolvent Spectral Encodings}
  \fancyfoot[C]{\thepage}
  \renewcommand{\headrulewidth}{0.4pt}
}

\titleformat{\section}{\Large\bfseries}{\thesection}{0.75em}{}
\titleformat{\subsection}{\large\bfseries}{\thesubsection}{0.65em}{}
\setlist[itemize]{leftmargin=1.55em,itemsep=0.25em,topsep=0.35em}
\setlist[enumerate]{leftmargin=1.75em,itemsep=0.25em,topsep=0.35em}
\setlength{\emergencystretch}{2em}

\newtheorem{theorem}{Theorem}[section]
\newtheorem{proposition}[theorem]{Proposition}
\newtheorem{lemma}[theorem]{Lemma}
\newtheorem{corollary}[theorem]{Corollary}
\theoremstyle{definition}
\newtheorem{definition}[theorem]{Definition}
\newtheorem{example}[theorem]{Example}
\theoremstyle{remark}
\newtheorem{remark}[theorem]{Remark}
\newtheorem{warning}[theorem]{Warning}

\newcommand{\R}{\mathbb R}
\newcommand{\C}{\mathbb C}
\newcommand{\Z}{\mathbb Z}
\newcommand{\N}{\mathbb N}
\newcommand{\T}{\mathbb T}
\newcommand{\Hcal}{\mathcal H}
\newcommand{\Scal}{\mathcal S}
\newcommand{\Tr}{\operatorname{Tr}}
\newcommand{\Spec}{\operatorname{Spec}}
\newcommand{\Dom}{\operatorname{Dom}}
\newcommand{\ran}{\operatorname{ran}}
\newcommand{\im}{\operatorname{im}}
\newcommand{\rank}{\operatorname{rank}}
\newcommand{\Vol}{\operatorname{Vol}}
\newcommand{\diag}{\operatorname{diag}}
\newcommand{\Id}{\mathrm I}
\newcommand{\e}{\mathrm e}
\newcommand{\ii}{\mathrm i}
\newcommand{\dd}{\,\mathrm d}
\newcommand{\norm}[1]{\left\lVert #1\right\rVert}
\newcommand{\abs}[1]{\left\lvert #1\right\rvert}
\newcommand{\set}[1]{\left\{#1\right\}}

\begin{document}
\thispagestyle{empty}

\begin{center}
{\LARGE\bfseries A Layered Framework for Compact-Resolvent\\[0.12cm]
Spectral Encodings\par}
\vspace{0.40cm}
{\large Stassis Stashkevichyus\par}
\vspace{0.08cm}
{\normalsize Stassis Research Program\par}
\vspace{0.05cm}
{\normalsize Independent Researcher, Vilnius, Lithuania\par}
\vspace{0.05cm}
{\normalsize ORCID:
\href{https://orcid.org/0009-0000-2294-705X}{0009-0000-2294-705X}\par}
\vspace{0.05cm}
{\normalsize Email:
\href{mailto:theobserver.of.multiverses@proton.me}{theobserver.of.multiverses@proton.me}\par}
\vspace{0.14cm}
{\normalsize Expository framework preprint -- revised 1 August 2026\par}
\end{center}

\begin{abstract}
Spectral generating functions are often discussed as though compactness of an
underlying carrier, discreteness of an operator spectrum, trace-class heat
evolution, and geometric realizability were interchangeable. They are not.
This expository framework separates these assertions into typed layers.
The operator layer consists of a nonnegative self-adjoint operator with compact
resolvent; heat-admissibility is imposed independently. Its spectral counting
measure
\[
  \mu_L=\sum_j m_j\delta_{\nu_j}
\]
produces the heat trace and its complex-time deformation as Laplace transforms
\[
  Z_L(z)=\int_{[0,\infty)}\e^{-z\lambda}\,\dd\mu_L(\lambda),
  \qquad \Re z>0.
\]
We prove the elementary holomorphy, differentiation, domination, and exact
imaginary-periodicity criteria for this transform; correct the factorization
criterion for quantities recoverable through a readout; and distinguish formal,
analytic, abstract diagonal, elliptic-geometric, equivariant, and physical
realization levels. Circle and flat-torus Laplacians provide normalized theta
examples. Two counterexamples show that a compact carrier need not produce
discrete spectrum and that compact resolvent need not produce a trace-class heat
semigroup. A rapidly growing abstract multiplicity model then exhibits the gap
between diagonal operator realizability and finite-dimensional elliptic
geometry. No new theorem in spectral geometry is claimed; the contribution is a
rigorous interface and failure-mode taxonomy assembled from standard operator
and heat-kernel theory.
\end{abstract}

\noindent\textbf{2020 Mathematics Subject Classification.}
Primary 47A10, 47A60; Secondary 58J50, 11F27.

\noindent\textbf{Keywords.}
Compact resolvent; heat-admissible operator; spectral counting measure;
complex-time heat trace; theta kernel; readout factorization; realization
hierarchy.

\vspace{0.25cm}
\noindent\textbf{Status and scope.}
This is an expository framework article. Except for elementary interface
propositions proved here for completeness, its analytic inputs are standard
results in functional analysis, elliptic theory, harmonic analysis, and compact
group representation theory. It does not propose a new spectral theory, a
physical compactification model, or an inverse spectral theorem.

\section{Introduction}

The phrase ``compact spectral geometry'' can conceal several logically
independent assertions. A carrier may be compact while a chosen operator has
continuous spectrum. A nonnegative self-adjoint operator may have compact
resolvent while its heat semigroup fails to be trace class for small positive
times. A formal power series may admit an abstract diagonal operator realization
while being incompatible with the Weyl and heat asymptotics of every
Laplace-type operator on a closed finite-dimensional manifold. Finally, a
representation label attached to an eigenspace is algebraic data, not by itself
a physical particle sector.

The purpose of this article is to make these boundaries explicit. The basic
chain is
\[
\boxed{
\text{operator hypotheses}
\longrightarrow
\text{spectral measure}
\longrightarrow
\text{analytic encoding}
\longrightarrow
\text{declared readout}.
}
\]
Each arrow has hypotheses. None may be replaced by an analogy.

\subsection{The layered contract}

We use the following hierarchy.

\begin{enumerate}
\item \textbf{Formal layer.} An expression
      \(\sum_i a_iq^{\lambda_i}\) is symbolic data until the index set,
      levels, coefficients, convergence regime, and branch convention are
      specified.
\item \textbf{Abstract operator layer.} A pair \((\Hcal,L)\), with \(L\)
      nonnegative self-adjoint and compact-resolvent, supplies a discrete
      spectral counting measure and finite spectral windows.
\item \textbf{Heat layer.} The extra condition
      \(\e^{-tL}\in\Scal_1\) for every \(t>0\) supplies a heat trace and a
      holomorphic complex-time trace.
\item \textbf{Geometric layer.} A closed manifold, Hermitian bundle, and
      Laplace-type operator supply compact resolvent, heat-admissibility, Weyl
      growth, and a local small-time expansion.
\item \textbf{Equivariant/readout layer.} A commuting compact-group action and a
      specified map identify symmetry labels and exactly which quantities
      descend through the map.
\item \textbf{Physical layer.} An action, units, observables, parameter
      identification, and empirical tests are additional data. They are outside
      the present framework.
\end{enumerate}

The separation between the abstract and geometric layers is essential. Every
suitable discrete multiplicity list has a diagonal realization, whereas local
elliptic realizations obey strong asymptotic constraints.

\subsection{Relation to standard frameworks}

Compact-resolvent spectral theory, trace ideals, heat kernels, and elliptic heat
asymptotics are classical; see
\cite{ReedSimonIV,Kato,SimonTrace,Davies,MinakPleijel,Gilkey1975,Seeley1967}.
The circle and torus formulas are direct instances of Fourier analysis and
Poisson summation \cite{SteinShakarchi,MumfordTheta}. The compact-group statements
are standard consequences of the spectral theorem and complete reducibility
\cite{Folland,FultonHarris}.

The datum used below is deliberately weaker than a spectral triple. It contains
no represented algebra, bounded-commutator condition, real structure, grading,
or dimension axiom, and therefore does not inherit conclusions of
noncommutative geometry \cite{Connes1994}. It also contains no spectral-action
functional \cite{ChamseddineConnes1997}. The present contribution is
organizational: it gives a minimal typed interface, records the exact descent
criterion for readouts, and isolates counterexamples that invalidate common
shortcuts.

\section{Compact resolvent and heat-admissibility}

\subsection{Operator data}

\begin{definition}[Abstract compact-resolvent datum]
An abstract compact-resolvent spectral datum is a pair \((\Hcal,L)\) such that
\(\Hcal\) is a complex separable Hilbert space and
\[
 L:\Dom(L)\subset\Hcal\longrightarrow\Hcal
\]
is densely defined, nonnegative, self-adjoint, and has compact resolvent.
Because \(L\geq0\), it is enough to require \((L+\Id)^{-1}\) to be compact.
\end{definition}

This definition does not mention a carrier. A carrier becomes relevant only
after a concrete realization of \(\Hcal\) and \(L\) has been specified.

\begin{theorem}[Standard compact-resolvent spectral theorem]
\label{thm:compact-resolvent}
Let \((\Hcal,L)\) be an abstract compact-resolvent datum.
\begin{enumerate}
\item If \(\dim\Hcal<\infty\), then \(L\) has a finite orthonormal eigenbasis.
\item If \(\dim\Hcal=\infty\), then there is an orthonormal basis
      \((e_j)_{j\geq0}\) and a sequence
      \[
        0\leq\lambda_0\leq\lambda_1\leq\cdots,\qquad
        \lambda_j\longrightarrow\infty,
      \]
      such that \(Le_j=\lambda_je_j\). Eigenvalues are repeated according to
      finite multiplicity.
\item For every bounded Borel set \(W\subset\R\), the spectral projection
      \(\mathbf 1_W(L)\) has finite rank.
\end{enumerate}
\end{theorem}

\begin{proof}
The compact self-adjoint operator \(R=(L+\Id)^{-1}\) has an orthonormal
eigenbasis for its nonzero spectrum together with a basis of \(\ker R\).
Here \(R\) is injective, hence \(\ker R=\{0\}\), so the nonzero eigenvectors
are complete. In infinite dimension, the nonzero eigenvalues of \(R\) have
finite multiplicity and can accumulate only at zero. If \(Re=\rho e\), then
\(e\in\ran R=\Dom(L)\) and \(Le=(\rho^{-1}-1)e\). This transfers the basis to
\(L\), gives finite multiplicities, and forces \(\lambda_j\to\infty\).
The finite-dimensional alternative must be stated separately. In either case,
a bounded set meets only finitely many eigenvalues counted with multiplicity,
so its spectral projection has finite rank. See
\cite[Chs.~VI--VII]{ReedSimonI} or \cite[Ch.~III]{Kato}.
\end{proof}

\begin{definition}[Spectral window]
For a bounded Borel set \(W\subset[0,\infty)\), the spectral window is
\[
  \Hcal_W:=\ran\mathbf 1_W(L).
\]
By Theorem~\ref{thm:compact-resolvent}, \(\Hcal_W\) is finite-dimensional.
\end{definition}

\subsection{Compact carrier is not an operator hypothesis}

\begin{example}[Compact carrier with continuous spectrum]
\label{ex:compact-carrier}
Let \(K=S^1\), \(\Hcal=L^2(S^1)\), and
\[
  (Lf)(\theta)=(2+\cos\theta)f(\theta).
\]
Then \(K\) is compact and \(L\) is bounded, nonnegative, and self-adjoint, but
\[
  \Spec(L)=[1,3],
\]
the essential range of \(2+\cos\theta\). Moreover,
\((L+\Id)^{-1}\) is multiplication by \((3+\cos\theta)^{-1}\), which is not
compact on \(L^2(S^1)\). Indeed, this multiplier has the bounded inverse given
by multiplication by \(3+\cos\theta\); if it were compact, the identity on the
infinite-dimensional space \(L^2(S^1)\) would be compact. Thus compactness of
the carrier alone gives neither compact resolvent nor discrete spectrum.
\end{example}

\begin{remark}
Compactness becomes effective only through an analytic theorem. For a
Laplace-type operator on a closed manifold, elliptic regularity maps the
resolvent into a higher Sobolev space and Rellich compactness returns it
compactly to \(L^2\). The conclusion comes from the operator and its domain,
not from the topology of the carrier in isolation.
\end{remark}

\subsection{Heat-admissibility is an independent condition}

\begin{definition}[Heat-admissibility]
A nonnegative self-adjoint operator \(L\) is heat-admissible if
\[
  \e^{-tL}\in\Scal_1(\Hcal)
  \qquad\text{for every }t>0,
\]
where \(\Scal_1\) is the trace-class ideal. Its heat trace is
\[
  Z_L(t):=\Tr(\e^{-tL}).
\]
\end{definition}

\begin{remark}
This is a naming convention for the present layer structure, not a new
summability notion. When \(L=D^2\), the same condition is commonly called
\(\theta\)-summability in spectral-triple terminology
\cite{Connes1994,SimonTrace}.
\end{remark}

\begin{example}[Compact resolvent without heat-admissibility]
\label{ex:log-spectrum}
On \(\ell^2(\N)\), define
\[
  Le_n=\log(n+1)e_n,\qquad
  \Dom(L)=
  \set{x:\sum_{n\geq1}\log^2(n+1)\abs{x_n}^2<\infty}.
\]
Then \(L\geq0\) is self-adjoint and
\((L+\Id)^{-1}\) is diagonal with entries tending to zero, hence compact.
Self-adjointness follows from the maximal-domain construction for a real
diagonal operator.
However,
\[
 \Tr(\e^{-tL})=\sum_{n\geq1}(n+1)^{-t},
\]
which diverges for \(0<t\leq1\). Compact resolvent therefore does not imply
heat-admissibility.
\end{example}

\begin{remark}[Growth criterion]
If \(\lambda_j\) are repeated with multiplicity, heat-admissibility is exactly
\[
  \sum_j\e^{-t\lambda_j}<\infty
  \quad\text{for every }t>0.
\]
It is a quantitative restriction on spectral growth, not merely a discreteness
condition.
\end{remark}

\section{Spectral measures and complex-time encodings}

\subsection{Counting measure and heat trace}

Let \(\nu_0<\nu_1<\cdots\) denote the distinct eigenvalues of a
compact-resolvent operator and let \(m_j\) be their multiplicities. The list is
finite in finite dimension; otherwise it is an infinite sequence with
\(\nu_j\to\infty\).

\begin{definition}[Spectral counting measure]
The spectral counting measure and counting function are
\[
  \mu_L:=\sum_jm_j\delta_{\nu_j},
  \qquad
  N_L(\Lambda):=\mu_L([0,\Lambda]).
\]
\end{definition}

The measure \(\mu_L\) is locally finite. If \(L\) is heat-admissible, then
\[
 Z_L(t)=\int_{[0,\infty)}\e^{-t\lambda}\,\dd\mu_L(\lambda)
       =\sum_jm_j\e^{-t\nu_j}.
\]
Thus the heat trace is the Laplace transform of a positive discrete measure.
This formulation keeps levels and multiplicities distinct and avoids assigning
spectral meaning to an undeclared coefficient sequence.

\subsection{Complex time}

\begin{definition}[Complex-time spectral trace]
For a heat-admissible \(L\) and \(z\in\C\) with \(\Re z>0\), define
\[
  Z_L(z):=\Tr(\e^{-zL})
  =\int_{[0,\infty)}\e^{-z\lambda}\,\dd\mu_L(\lambda).
\]
Writing \(z=t-\ii\phi\), with \(t>0\) and \(\phi\in\R\), gives
\[
  Z_L(t-\ii\phi)
  =\sum_jm_j\e^{-t\nu_j}\e^{\ii\phi\nu_j}.
\]
The phase parameter is initially real, not an element of
\(\R/2\pi\Z\).
\end{definition}

\begin{proposition}[Holomorphic trace and derivative bounds]
\label{prop:holomorphic}
Let \(L\) be heat-admissible. Then:
\begin{enumerate}
\item \(\e^{-zL}\) is trace class for every \(\Re z>0\), and
      \[
        \norm{\e^{-zL}}_1=Z_L(\Re z).
      \]
\item \(Z_L\) is holomorphic on the right half-plane and, for every
      \(r\in\N_0\),
      \[
        Z_L^{(r)}(z)
        =(-1)^r\int\lambda^r\e^{-z\lambda}\,\dd\mu_L(\lambda).
      \]
\item The domination estimate
      \[
        \abs{Z_L(z)}\leq Z_L(\Re z)
      \]
      holds.
\end{enumerate}
\end{proposition}

\begin{proof}
Functional calculus gives
\[
 \e^{-zL}=\e^{-(\Re z)L}\e^{-\ii(\Im z)L}.
\]
The second factor is unitary and commutes with the positive trace-class first
factor, giving the trace norm identity. On every compact subset of
\(\{\Re z>0\}\), there is \(a>0\) with \(\Re z\geq a\). The defining series is
dominated by \(\sum_jm_j\e^{-a\nu_j}<\infty\), so it converges locally uniformly.
For derivatives, use
\[
 \lambda^r\e^{-a\lambda}
 \leq C_{r,a}\e^{-a\lambda/2}
\]
and heat-admissibility at \(a/2\). Termwise differentiation is therefore
valid. The final estimate follows from the triangle inequality.
\end{proof}

\begin{proposition}[Exact imaginary-periodicity criterion]
\label{prop:periodicity}
Let \(L\) be heat-admissible and let \(T>0\). The following are equivalent:
\begin{enumerate}
\item \(Z_L(z+\ii T)=Z_L(z)\) for all \(\Re z>0\);
\item \(T\nu_j\in2\pi\Z\) for every eigenvalue \(\nu_j\).
\end{enumerate}
In particular, \(2\pi\)-periodicity in \(\phi\) is automatic only when all
levels are integers.
\end{proposition}

\begin{proof}
The second condition immediately implies the first. Conversely, evaluate the
assumed identity at real \(z=t>0\):
\[
 \sum_jm_j\bigl(\e^{-\ii T\nu_j}-1\bigr)\e^{-t\nu_j}=0.
\]
Order the distinct eigenvalues increasingly and set
\(c_j=m_j(\e^{-\ii T\nu_j}-1)\). For any fixed \(t_0>0\) and \(t\geq t_0\),
\[
 \abs{c_j}\e^{-t(\nu_j-\nu_0)}
 \leq 2\e^{t_0\nu_0}m_j\e^{-t_0\nu_j},
\]
and the majorant is summable by heat-admissibility. Multiplying the identity
by \(\e^{t\nu_0}\) and applying dominated convergence as \(t\to\infty\)
therefore gives \(c_0=0\). Remove that term and apply the same argument to
the tail beginning at \(\nu_1\), then iterate. Hence
\(\e^{-\ii T\nu_j}=1\), equivalently \(T\nu_j\in2\pi\Z\), for every \(j\).
\end{proof}

\begin{remark}[Safe \(q\)-notation]
If the levels are nonnegative integers, one may set \(q=\e^{-z}\), \(\abs q<1\),
and write \(Z_L(z)=\sum_nm_nq^n\). For general real levels, \(q^\lambda\)
requires a choice of logarithm. The invariant notation is
\(\e^{-z\lambda}\).
\end{remark}

\subsection{Mellin transform and the positive spectral zeta function}

\begin{definition}
Let \(h_L=\dim\ker L\). The positive spectral zeta function is
\[
 \zeta_L(s):=\int_{(0,\infty)}\lambda^{-s}\,\dd\mu_L(\lambda)
            =\sum_{\nu_j>0}m_j\nu_j^{-s}
\]
in its half-plane of absolute convergence.
\end{definition}

\begin{proposition}[Mellin identity with explicit hypothesis]
\label{prop:mellin}
Let \(L\) be heat-admissible, let \(s\in\C\), let
\(\sigma=\Re s>0\), and assume
\(\sum_{\nu_j>0}m_j\nu_j^{-\sigma}<\infty\). Then
\[
  \zeta_L(s)
  =\frac{1}{\Gamma(s)}
    \int_0^\infty t^{s-1}\bigl(Z_L(t)-h_L\bigr)\,\dd t.
\]
\end{proposition}

\begin{proof}
For \(\lambda>0\),
\[
 \lambda^{-s}=\frac1{\Gamma(s)}
 \int_0^\infty t^{s-1}\e^{-t\lambda}\,\dd t.
\]
Moreover, Tonelli's theorem gives the exact absolute-integrability test
\[
 \int_0^\infty
 \sum_{\nu_j>0}m_jt^{\sigma-1}\e^{-t\nu_j}\,\dd t
 =
 \Gamma(\sigma)\sum_{\nu_j>0}m_j\nu_j^{-\sigma}<\infty.
\]
Fubini's theorem therefore gives the identity. No analytic continuation is
asserted.
\end{proof}

\section{A realization hierarchy}

\subsection{Formal and abstract realizations}

\begin{definition}[Level encoding]
A level encoding is data \((I,\lambda,a)\), where \(I\) is countable,
\(\lambda:I\to[0,\infty)\), and \(a:I\to\C\). It is \emph{locally finite} if
\[
 \#\set{i\in I:a_i\neq0,\ \lambda_i\leq\Lambda}<\infty
 \qquad(\Lambda<\infty).
\]
On an open set \(D\subset\C\), it defines an analytic encoding if the series
\[
 F(z)=\sum_{i\in I}a_i\e^{-z\lambda_i}
\]
converges normally on \(D\), meaning that for every compact \(K\subset D\),
\[
 \sum_{i\in I}\sup_{z\in K}
 \abs{a_i\e^{-z\lambda_i}}<\infty.
\]
Normal convergence makes \(F\) holomorphic. Under local finiteness, the grouped
coefficient at a level \(\nu\) is the finite sum
\[
 b_\nu:=\sum_{\lambda_i=\nu}a_i.
\]
The encoding can be an abstract heat trace only if every \(b_\nu\) is a
nonnegative integer, the support \(\set{\nu:b_\nu\neq0}\) is locally finite in
\([0,\infty)\), and
\(\sum_\nu b_\nu\e^{-t\nu}<\infty\) for every \(t>0\).
\end{definition}

\begin{proposition}[Canonical diagonal realization]
\label{prop:diagonal}
Let
\[
 0\leq\lambda_0<\lambda_1<\cdots,\qquad
 \lambda_n\to\infty,\qquad m_n\in\N_{>0}.
\]
On
\[
 \Hcal=\bigoplus_{n\geq0}\C^{m_n}
\]
define \(L\) to act by \(\lambda_n\Id\) on the \(n\)-th summand, with domain
\[
 \Dom(L)=
 \set{x=(x_n):\sum_{n\geq0}\lambda_n^2\norm{x_n}^2<\infty}.
\]
Then \(L\geq0\) is self-adjoint with compact resolvent.
It is heat-admissible if and only if
\[
 \sum_{n\geq0}m_n\e^{-t\lambda_n}<\infty
 \quad\text{for every }t>0.
\]
The resulting pure-point operator is unique up to unitary equivalence among
operators with exactly these eigenvalues and multiplicities.
For a finite level list, the same statement holds on a finite-dimensional
direct sum, with the condition \(\lambda_n\to\infty\) omitted.
\end{proposition}

\begin{proof}
The displayed domain is the maximal domain of this real diagonal operator, so
\(L\) is self-adjoint and nonnegative.
The resolvent is block diagonal with block norm
\((1+\lambda_n)^{-1}\to0\), hence it is the norm limit of finite-rank
truncations. The trace-class criterion follows by summing the positive
eigenvalues \(m_n\e^{-t\lambda_n}\) of \(\e^{-tL}\). Unitary equivalence follows
by choosing unitary identifications between corresponding finite-dimensional
eigenspaces.
\end{proof}

\begin{warning}
Proposition~\ref{prop:diagonal} makes abstract realizability inexpensive. It does not
produce a carrier, locality, differential expression, principal symbol,
dimension, metric, boundary condition, or physical interpretation.
\end{warning}

\subsection{Elliptic geometric realization}

\begin{definition}[Closed Laplace-type realization]
A closed Laplace-type realization consists of a closed smooth
\(d\)-dimensional Riemannian manifold \((M,g)\), \(d\geq1\), a finite-rank
Hermitian bundle
\(E\to M\), and the nonnegative self-adjoint realization on \(L^2(M,E)\) of a
second-order differential operator \(L\) whose principal symbol is
\[
 \sigma_2(L)(x,\xi)=\abs{\xi}_g^2\Id_{E_x}.
\]
\end{definition}

\begin{theorem}[Standard elliptic package]
\label{thm:elliptic}
For a closed Laplace-type realization:
\begin{enumerate}
\item \(L\) has compact resolvent and is heat-admissible;
\item its counting function satisfies the Weyl law
      \[
        N_L(\Lambda)
        \sim
        \frac{\rank(E)\,\omega_d}{(2\pi)^d}\Vol(M)\Lambda^{d/2},
        \qquad \Lambda\to\infty;
      \]
\item its heat trace has a complete small-time expansion
      \[
        Z_L(t)
        \sim
        (4\pi t)^{-d/2}
        \sum_{k=0}^\infty A_k(L)t^k,
        \qquad t\downarrow0,
      \]
      with \(A_0(L)=\rank(E)\Vol(M)\) and higher coefficients determined by
      local geometric and bundle data.
\end{enumerate}
\end{theorem}

\begin{proof}[Source]
Compact resolvent follows from elliptic regularity and Rellich compactness.
The heat-kernel construction, local expansion, and Weyl law are classical;
see \cite{MinakPleijel,Gilkey1975,Davies,BGV,Chavel,Seeley1967}. This theorem
is quoted as a standard package; it is not presented as a new result.
\end{proof}

\begin{remark}
The displayed asymptotics are necessary tests for a closed Laplace-type
realization. They are not sufficient: matching a finite or even an asymptotic
list of coefficients does not construct a manifold and operator.
Boundaries, singular spaces, noncompact carriers, nonlocal operators, and
elliptic operators of order other than two have different theorem packages and
must be declared separately.
\end{remark}

\begin{example}[Abstract heat trace excluded by elliptic asymptotics]
\label{ex:fast-multiplicity}
Let \(\lambda_n=n\) and \(m_n=\lceil\e^{\sqrt n}\rceil\) for \(n\geq1\).
Proposition~\ref{prop:diagonal} gives a compact-resolvent operator. It is
heat-admissible because
\(\sum_n\e^{\sqrt n-tn}<\infty\) for every \(t>0\).
For \(0<t\leq1/4\), choose
\(n_t=\lfloor(4t^2)^{-1}\rfloor\). Then, for a constant \(c>0\),
\[
 Z_L(t)\geq m_{n_t}\e^{-tn_t}\geq \e^{c/t}
\]
once \(t\) is sufficiently small. Hence \(Z_L(t)\) grows faster than every
power of \(t^{-1}\). It cannot be the heat trace of a Laplace-type operator on
any closed finite-dimensional manifold, because that would contradict
Theorem~\ref{thm:elliptic}.
\end{example}

\begin{proof}
For fixed \(t>0\), the inequality \(\sqrt n\leq tn/2\) holds once
\(n\geq4t^{-2}\); hence the tail of
\(\sum_n\e^{\sqrt n-tn}\) is bounded by a geometric series
\(\sum_n\e^{-tn/2}\). This proves heat-admissibility.

For \(t\leq1/4\), the number \(x=(4t^2)^{-1}\) is at least \(4\), so
\(n_t\geq x/2=(8t^2)^{-1}\), while \(tn_t\leq1/(4t)\). Therefore
\[
 \sqrt{n_t}-tn_t
 \geq
 \left(\frac1{2\sqrt2}-\frac14\right)\frac1t.
\]
The coefficient in parentheses is positive, which proves the bound.
\end{proof}

\subsection{Layer table}

\begin{center}
\small
\begin{tabularx}{\textwidth}{@{}>{\raggedright\arraybackslash}p{0.19\textwidth}
>{\raggedright\arraybackslash}p{0.29\textwidth}
>{\raggedright\arraybackslash}X@{}}
\toprule
Layer & Required data & What is justified \\
\midrule
Formal
& \(I,\lambda,a\), syntax
& A symbolic level expression only. \\
Convergent analytic
& Normal convergence on an open set
& A holomorphic function on that set. \\
Abstract spectral
& Integer multiplicities and a self-adjoint realization
& A pure-point spectrum; compact resolvent if levels escape to infinity. \\
Heat
& \(\e^{-tL}\in\Scal_1\) for all \(t>0\)
& Trace, spectral Laplace transform, complex-time encoding. \\
Elliptic geometric
& \(M,E,L\), domain and ellipticity
& Weyl law, local heat expansion, geometric coefficients. \\
Equivariant
& Strongly continuous unitary action commuting with \(L\)
& Representation structure on finite eigenspaces. \\
Physical
& Dynamics, units, observables, calibration and tests
& Model-dependent physical statements only. \\
\bottomrule
\end{tabularx}
\end{center}

\section{Circle and flat-torus models}

We use the Fourier-transform convention
\[
 \widehat f(\xi)=\int_{\R^d}f(x)\e^{-2\pi\ii x\cdot\xi}\,\dd x,
\]
for which the lattice Poisson formula is
\(\sum_{n\in\Z^d}f(n)=\sum_{k\in\Z^d}\widehat f(k)\) on the Schwartz class.

\subsection{Circle}

Let \(S^1=\R/2\pi\Z\) and equip it with normalized Haar measure
\(\dd\theta/(2\pi)\). For \(R>0\), define
\[
  L_R=-R^{-2}\frac{\dd^2}{\dd\theta^2},
  \qquad \Dom(L_R)=H^2_{\mathrm{per}}(S^1).
\]

\begin{proposition}[Circle spectrum and theta trace]
The modes \(e_n(\theta)=\e^{\ii n\theta}\), \(n\in\Z\), form an orthonormal
eigenbasis with
\[
  L_Re_n=\frac{n^2}{R^2}e_n.
\]
Consequently,
\[
  Z_{L_R}(t)=\sum_{n\in\Z}\e^{-tn^2/R^2}
  =
  R\sqrt{\frac{\pi}{t}}
  \sum_{k\in\Z}\e^{-\pi^2R^2k^2/t}.
\]
\end{proposition}

\begin{proof}
The spectral statement is Fourier theory on \(S^1\). Apply Poisson summation to
\(x\mapsto\e^{-tx^2/R^2}\) to obtain the second expression.
\end{proof}

The equality relates short and long heat scales, but it is a precise identity for
this quadratic spectrum, not a universal duality principle.

\subsection{Flat tori}

Let \(\T^d=(\R/2\pi\Z)^d\) with normalized Haar measure, let \(A\) be a real
symmetric positive-definite \(d\times d\) matrix, and define
\[
 L_A=-\sum_{a,b=1}^d A_{ab}\partial_a\partial_b
 \quad\text{on}\quad H^2(\T^d).
\]

\begin{proposition}[Flat-torus spectrum and Poisson transform]
\label{prop:torus}
The normalized Fourier modes
\[
 e_n(\theta)=\e^{\ii n\cdot\theta},\qquad n\in\Z^d,
\]
satisfy
\[
 L_Ae_n=(n^{\mathsf T}An)e_n.
\]
The heat trace is
\[
 Z_A(t)=\sum_{n\in\Z^d}\e^{-t n^{\mathsf T}An},
\]
and Poisson summation gives
\[
 Z_A(t)=
 \frac{\pi^{d/2}}{t^{d/2}\sqrt{\det A}}
 \sum_{k\in\Z^d}
 \exp\left(-\frac{\pi^2}{t}k^{\mathsf T}A^{-1}k\right).
\]
\end{proposition}

\begin{proof}
Differentiating \(e_n\) gives the eigenvalue \(n^{\mathsf T}An\).
Positive definiteness yields \(n^{\mathsf T}An\geq c\abs n^2\), so the heat
trace converges. The Fourier transform of
\(x\mapsto\e^{-t x^{\mathsf T}Ax}\) is
\[
 \xi\longmapsto
 \frac{\pi^{d/2}}{t^{d/2}\sqrt{\det A}}
 \exp\left(-\frac{\pi^2}{t}\xi^{\mathsf T}A^{-1}\xi\right),
\]
and the lattice Poisson formula proves the identity.
\end{proof}

\begin{remark}[Degeneracy]
The multiplicity of \(\lambda\) is the finite representation number
\[
 r_A(\lambda)=
 \#\set{n\in\Z^d:n^{\mathsf T}An=\lambda}.
\]
A statement about one Fourier label \(n\) must not be confused with a statement
about the full eigenspace when several lattice vectors have the same quadratic
value.
\end{remark}

\subsection{Tensor sums}

\begin{proposition}[Product heat trace]
Let \(L_i\) be a nonnegative heat-admissible self-adjoint operator on
\(\Hcal_i\), \(i=1,2\). On \(\Hcal_1\widehat\otimes\Hcal_2\), let \(L\) be the
nonnegative self-adjoint operator associated with the closed quadratic form
\[
 \Dom(q)=
 \Dom(L_1^{1/2}\otimes\Id)\cap
 \Dom(\Id\otimes L_2^{1/2}),
\qquad
 q[u]=
 \norm{(L_1^{1/2}\otimes\Id)u}^2+
 \norm{(\Id\otimes L_2^{1/2})u}^2
\]
for \(u\in\Dom(q)\). Equivalently, \(L\) is the self-adjoint tensor sum
\(L_1\otimes\Id+\Id\otimes L_2\).
Then \(L\) is heat-admissible and, for \(\Re z>0\),
\[
 \e^{-zL}=\e^{-zL_1}\otimes\e^{-zL_2},
 \qquad
 Z_L(z)=Z_{L_1}(z)Z_{L_2}(z).
\]
\end{proposition}

\begin{proof}
The spectral measures of the two tensor factors commute, and their joint
functional calculus represents \(L\) by the function
\((\lambda_1,\lambda_2)\mapsto\lambda_1+\lambda_2\). This gives the
exponential identity. For real \(z=t>0\), the tensor product of the two
positive trace-class heat operators is trace class, proving
heat-admissibility. The same trace-ideal fact for complex \(z\), together with
\(\Tr(A\otimes B)=\Tr(A)\Tr(B)\), gives the trace formula.
\end{proof}

\section{Readouts and equivariant labels}

\subsection{Exact factorization through a readout}

\begin{definition}[Readout and recoverability]
Let \(\Pi:X\to Y\) be a specified map and \(Q:X\to Z\) a quantity. We say that
\(Q\) is recoverable from \(\Pi\) if it factors through the actual image
\(\Pi(X)\).
\end{definition}

\begin{proposition}[Correct fiber criterion]
\label{prop:fiber}
There exists a unique map
\[
 \widehat Q:\Pi(X)\longrightarrow Z
\quad\text{such that}\quad
 Q=\widehat Q\circ\Pi
\]
if and only if
\[
 \Pi(x)=\Pi(x')\Longrightarrow Q(x)=Q(x')
 \qquad(x,x'\in X).
\]
If \(\Pi\) is surjective, then \(\Pi(X)=Y\). Without surjectivity, the criterion
does not determine values of an extension \(Y\to Z\) outside \(\Pi(X)\).
\end{proposition}

\begin{proof}
Factorization implies fiber constancy. Conversely, define
\(\widehat Q(\Pi(x))=Q(x)\). Fiber constancy makes the definition independent
of the representative and uniqueness holds on \(\Pi(X)\).
\end{proof}

\begin{corollary}[Continuous descent]
Let \(X,Y,Z\) be topological spaces and let \(\Pi:X\to Y\) be a surjective
quotient map. A continuous map \(Q:X\to Z\) admits a unique continuous
factorization \(Q=\widehat Q\circ\Pi\) if and only if it is constant on the
fibers of \(\Pi\).
\end{corollary}

\begin{proof}
Apply Proposition~\ref{prop:fiber}. The defining property of the quotient topology says
that \(\widehat Q\) is continuous exactly when
\(\widehat Q\circ\Pi=Q\) is continuous.
\end{proof}

\begin{remark}
If smooth, measurable, Lipschitz, or other regular descent is required, the
corresponding category needs its own hypotheses. Set-theoretic fiber constancy
alone proves only set-theoretic factorization.
\end{remark}

\subsection{Compact-group symmetry}

\begin{definition}[Equivariant spectral datum]
Let \(G\) be a compact group with a strongly continuous unitary representation
\(U:G\to\mathcal U(\Hcal)\). The datum \((\Hcal,L,U)\) is equivariant if
\[
 U_g\Dom(L)=\Dom(L),
 \qquad
 LU_gx=U_gLx
 \quad(x\in\Dom(L)).
\]
\end{definition}

\begin{proposition}[Static symmetry labels]
\label{prop:static-symmetry}
If \(L\) has compact resolvent and \((\Hcal,L,U)\) is equivariant, then every
spectral projection \(\mathbf1_W(L)\) commutes with \(U_g\). In particular,
each finite-dimensional eigenspace
\[
 E_\lambda=\ker(L-\lambda\Id)
\]
is a unitary \(G\)-representation and decomposes into irreducibles.
\end{proposition}

\begin{proof}
The domain invariance and equality \(LU_g=U_gL\) imply
\[
 U_g(L+\Id)^{-1}=(L+\Id)^{-1}U_g.
\]
The bounded functional calculus therefore gives commutation with every bounded
Borel function of \(L\), hence with each spectral projection. The eigenspaces
are finite-dimensional by compact resolvent. Their restricted \(G\)-actions
are unitary, so invariant subspaces have invariant orthogonal complements;
iteration gives a finite direct sum of irreducibles.
\end{proof}

\begin{warning}
Proposition~\ref{prop:static-symmetry} is a static statement. It does not assert that the
multiplicity of a fixed eigenvalue remains constant under perturbation.
Eigenvalues may split, merge, or cross. Quantitative stability of an isolated
cluster requires a declared spectral gap and control of the corresponding Riesz
projection; that problem is not treated here.
\end{warning}

\subsection{A corrected toric internal model}

Let \(V\) be a finite-dimensional unitary \(G\)-representation and consider
\[
 \Hcal=L^2(\T^d)\otimes V,
 \qquad
 L=L_A\otimes\Id_V.
\]
Here \(G\) acts trivially on \(L^2(\T^d)\) and through the given
representation on \(V\).

\begin{proposition}[Full eigenspace at a toric level]
\label{prop:toric-internal}
For every eigenvalue \(\lambda\) of \(L_A\),
\[
 E_\lambda(L)
 =
 \bigoplus_{\substack{n\in\Z^d\\n^{\mathsf T}An=\lambda}}
 \C e_n\otimes V.
\]
Consequently, as a \(G\)-representation,
\[
 E_\lambda(L)\cong V^{\oplus r_A(\lambda)}.
\]
If \(H\subset G\) is closed and
\[
 V\!\downarrow_H\cong\bigoplus_\alpha b_\alpha W_\alpha,
\]
then
\[
 E_\lambda(L)\!\downarrow_H
 \cong
 \bigoplus_\alpha
 \bigl(r_A(\lambda)b_\alpha\bigr)W_\alpha.
\]
\end{proposition}

\begin{proof}
Proposition~\ref{prop:torus} identifies all Fourier modes with eigenvalue \(\lambda\).
Tensoring each Fourier line with \(V\) gives the full eigenspace. The group acts
trivially on the Fourier multiplicity space and through the given
representation on \(V\), so multiplicities multiply by \(r_A(\lambda)\).
\end{proof}

\begin{remark}
This proposition supplies representation-theoretic labels attached to a
spectral level. It does not identify those labels with physical fields,
charges, masses, or interactions.
\end{remark}

\section{Claim boundary and admissibility test}

The framework is intentionally narrow. The following table records the main
logical boundaries.

\begin{center}
\small
\begin{tabularx}{\textwidth}{@{}>{\raggedright\arraybackslash}p{0.37\textwidth}
>{\raggedright\arraybackslash}X@{}}
\toprule
Invalid shortcut & Missing hypothesis or conclusion \\
\midrule
Compact carrier \(\Rightarrow\) discrete spectrum
& False by Example~\ref{ex:compact-carrier}; an operator theorem such as ellipticity
  and compact Sobolev embedding is required. \\
Compact resolvent \(\Rightarrow\) heat trace for every \(t>0\)
& False by Example~\ref{ex:log-spectrum}; trace-class summability is separate. \\
Formal \(q\)-series \(\Rightarrow\) spectrum
& Levels, integer multiplicities, convergence, and an operator realization are
  missing. \\
Abstract diagonal realization \(\Rightarrow\) compact geometry
& False in general; Weyl and heat asymptotics impose geometric obstructions,
  as in Example~\ref{ex:fast-multiplicity}. \\
\(\phi\in\R/2\pi\Z\) for arbitrary real levels
& False unless the spectrum satisfies the arithmetic condition in
  Proposition~\ref{prop:periodicity}. \\
Fiber constancy \(\Rightarrow\) a map on all of \(Y\)
& It canonically produces a map only on \(\Pi(X)\); surjectivity or a separate
  extension is needed. \\
Equivariant eigenspace \(\Rightarrow\) perturbative label stability
& False without an isolated cluster, a gap, and an equivariant perturbation
  theorem. \\
Representation branch \(\Rightarrow\) physical sector
& Dynamics, units, observables, and empirical identification remain absent. \\
\bottomrule
\end{tabularx}
\end{center}

\subsection{Audit protocol}

For use as a diagnostic rather than as a theorem, a proposed spectral encoding
should state, in order:
\begin{enumerate}
\item the index set and level function;
\item the meaning and sign/integrality of the coefficients;
\item the convergence domain;
\item the operator and its self-adjoint domain, if an operator claim is made;
\item compact-resolvent and heat-admissibility hypotheses separately;
\item the carrier, bundle, locality, and ellipticity data, if a geometric claim
      is made;
\item the group action and commutation relations, if symmetry labels are used;
\item the readout map and the category in which descent is asserted;
\item any additional bridge data needed for a physical interpretation.
\end{enumerate}

Items 1--3 define the analytic expression. Items 4--5 justify abstract spectral
and heat language. Item 6 is needed to pass from an abstract operator to
geometry. Items 7--8 type symmetry and recoverability statements. Item 9 is
logically external to the mathematical spectral package.

\section{Conclusion}

The correct foundational implication is
\[
\boxed{
L=L^*\geq0,\quad (L+\Id)^{-1}\text{ compact}
\quad\Longrightarrow\quad
\text{discrete finite-multiplicity spectrum},
}
\]
with
\[
\boxed{
\e^{-tL}\in\Scal_1\ \text{for every }t>0
\quad\Longrightarrow\quad
Z_L(z)=\int\e^{-z\lambda}\,\dd\mu_L(\lambda),\quad\Re z>0.
}
\]
Compactness of a carrier is not a substitute for the first hypothesis, and
compact resolvent is not a substitute for the second. Closed Laplace-type
geometry supplies both and adds Weyl and local heat asymptotics. Abstract
diagonal realizability supplies neither locality nor geometry.

The complex-time formulation removes the ambiguous use of a periodic phase:
periodicity is a theorem only when the spectrum lies in the corresponding
arithmetic lattice. The corrected readout criterion likewise produces a
canonical factor only on the actual image unless surjectivity is imposed.
Finally, equivariance gives representation structure on eigenspaces but not
perturbative stability or physical interpretation.

These distinctions make the framework suitable as an expository companion to
more focused research. Natural research questions -- geometric realizability
obstructions, quantitative equivariant cluster stability, and branch-resolved
heat-trace inversion -- require separate theorems and are intentionally not
claimed here.

\begin{thebibliography}{99}

\bibitem{BGV}
N.~Berline, E.~Getzler, and M.~Vergne,
\emph{Heat Kernels and Dirac Operators},
Grundlehren der mathematischen Wissenschaften 298,
Springer-Verlag, Berlin, 1992.
\href{https://doi.org/10.1007/978-3-642-58088-8}
{doi:10.1007/978-3-642-58088-8}.

\bibitem{ChamseddineConnes1997}
A.~H.~Chamseddine and A.~Connes,
The spectral action principle,
\emph{Communications in Mathematical Physics} \textbf{186} (1997), 731--750.
\href{https://doi.org/10.1007/s002200050126}{doi:10.1007/s002200050126}.

\bibitem{Chavel}
I.~Chavel,
\emph{Eigenvalues in Riemannian Geometry},
Pure and Applied Mathematics 115,
Academic Press, Orlando, 1984.

\bibitem{Connes1994}
A.~Connes,
\emph{Noncommutative Geometry},
Academic Press, San Diego, 1994.

\bibitem{Davies}
E.~B.~Davies,
\emph{Heat Kernels and Spectral Theory},
Cambridge Tracts in Mathematics 92,
Cambridge University Press, Cambridge, 1989.
\href{https://doi.org/10.1017/CBO9780511566158}{doi:10.1017/CBO9780511566158}.

\bibitem{Folland}
G.~B.~Folland,
\emph{A Course in Abstract Harmonic Analysis},
2nd ed., CRC Press, Boca Raton, 2016.
\href{https://doi.org/10.1201/b19172}{doi:10.1201/b19172}.

\bibitem{FultonHarris}
W.~Fulton and J.~Harris,
\emph{Representation Theory: A First Course},
Graduate Texts in Mathematics 129, Springer, New York, 1991.
\href{https://doi.org/10.1007/978-1-4612-0979-9}
{doi:10.1007/978-1-4612-0979-9}.

\bibitem{Gilkey1975}
P.~B.~Gilkey,
The spectral geometry of a Riemannian manifold,
\emph{Journal of Differential Geometry} \textbf{10} (1975), 601--618.
\href{https://doi.org/10.4310/jdg/1214433164}{doi:10.4310/jdg/1214433164}.

\bibitem{Kato}
T.~Kato,
\emph{Perturbation Theory for Linear Operators},
2nd ed., Springer-Verlag, Berlin, 1976; reprinted in Classics in Mathematics,
Springer, 1995.
\href{https://doi.org/10.1007/978-3-642-66282-9}
{doi:10.1007/978-3-642-66282-9}.

\bibitem{MinakPleijel}
S.~Minakshisundaram and {\AA}.~Pleijel,
Some properties of the eigenfunctions of the Laplace-operator on Riemannian
manifolds,
\emph{Canadian Journal of Mathematics} \textbf{1} (1949), 242--256.
\href{https://doi.org/10.4153/CJM-1949-021-5}{doi:10.4153/CJM-1949-021-5}.

\bibitem{MumfordTheta}
D.~Mumford,
\emph{Tata Lectures on Theta I},
Progress in Mathematics 28, Birkh\"auser, Boston, 1983.

\bibitem{ReedSimonI}
M.~Reed and B.~Simon,
\emph{Methods of Modern Mathematical Physics I: Functional Analysis},
rev. ed., Academic Press, San Diego, 1980.

\bibitem{ReedSimonIV}
M.~Reed and B.~Simon,
\emph{Methods of Modern Mathematical Physics IV: Analysis of Operators},
Academic Press, New York, 1978.

\bibitem{Seeley1967}
R.~T.~Seeley,
Complex powers of an elliptic operator,
in \emph{Singular Integrals}, Proceedings of Symposia in Pure Mathematics 10,
American Mathematical Society, Providence, 1967, pp.~288--307,
MR~0237943.

\bibitem{SimonTrace}
B.~Simon,
\emph{Trace Ideals and Their Applications},
2nd ed., Mathematical Surveys and Monographs 120,
American Mathematical Society, Providence, 2005.
\href{https://doi.org/10.1090/surv/120}{doi:10.1090/surv/120}.

\bibitem{SteinShakarchi}
E.~M.~Stein and R.~Shakarchi,
\emph{Fourier Analysis: An Introduction},
Princeton Lectures in Analysis I,
Princeton University Press, Princeton, 2003.

\end{thebibliography}

\end{document}
